% Source pp. 61-72 (Lectures7-12.pdf).
\section{Interaction. \texorpdfstring{$\varphi^4$}{phi**4} theory}

The KG theory describes free particles. Its action is quadratic in the field
variables. We want to add interactions. This is done by adding new terms to the
action. It should be done in a way that preserves the locality. We can modify
the Lagrangian density by adding local terms involving higher powers of the
field(s). The simplest example is known as the $\varphi^4$ theory, which we
consider next.

The $\varphi^4$ theory is described by the Euclidean action
\[
\mathcal A[\varphi]=\int \dd^4x\,\left[\frac12\,(\partial_\mu\varphi)^2+\frac{m^2}{2}\,\varphi^2+\frac{\lambda}{4!}\,\varphi^4\right]\,.
\]
Here $\lambda$ is a new parameter, the ``coupling constant''.

As in the KG theory, one can be interested in the correlation functions
\[
\bigl\langle\varphi(x_1)\cdots\varphi(x_N)\bigr\rangle=Z^{-1}\int\bigl[D\varphi\bigr]\,\varphi(x_1)\cdots\varphi(x_N)\,e^{-\mathcal A[\varphi]}\,,\qquad
Z=\int\bigl[D\varphi\bigr]\,e^{-\mathcal A[\varphi]}\,.
\]
Now the action is not quadratic in the integration variables $\varphi$, and
the integral cannot be evaluated in an explicit form.

\begin{exercise}
Performing the change of variables $\varphi(x)\to\varphi(x)+\varepsilon(x)$
in the functional integral defining the $N$-point correlation functions in
$\varphi^4$ theory, derive the equation
\begin{multline*}
\bigl(m^2-\partial_\mu^2\bigr)\bigl\langle\varphi(x)\,\varphi(x_1)\cdots\varphi(x_N)\bigr\rangle
+\frac{\lambda}{3!}\,\bigl\langle\varphi^3(x)\,\varphi(x_1)\cdots\varphi(x_N)\bigr\rangle\\
=\sum_{i=1}^N\delta^{(4)}(x-x_i)\,\bigl\langle\varphi(x_1)\cdots\widehat{\varphi(x_i)}\cdots\varphi(x_N)\bigr\rangle\,.
\end{multline*}
Equations of this kind are generally known as the \emph{Schwinger-Dyson
equations}.
\end{exercise}

Unlike the KG theory, the Schwinger-Dyson equation in the interacting theory is
not closed equation for the correlation function $\langle\varphi(x)\dots\rangle$,
since it involves $\langle\varphi^3(x)\dots\rangle$ which has no useful
expression in terms of the function $\langle\varphi(x)\dots\rangle$ itself.
Even the definition of the composite field $\varphi^3$ may be a non-trivial
problem (we are going to consider it later on).

Practical way of computations is the \emph{perturbation theory}, which
generates a power series expansion in the coupling parameter $\lambda$.

\subsection*{Perturbation Theory}

Perturbation theory allows one to evaluate the correlation functions as the
power series in the coupling parameter $\lambda$. One splits the action into
the free-field part (quadratic in the field variables $\varphi$), and the
interaction part, which includes the terms of higher order:
\[
\mathcal A=\mathcal A_0+\mathcal A_I\,.
\]
The exponential $\exp(-\mathcal A)$ is then expanded as
\[
e^{-\mathcal A}=e^{-\mathcal A_0}\left(1-\mathcal A_I+\frac12\,\mathcal A_I^2+\dots\right)\,.
\]
Since $\mathcal A_0$ is quadratic, each term can be evaluated using the Wick's
theorem.

To see how it works, consider the $\varphi^4$ theory. In this case the natural
way of splitting is
\[
\mathcal A_0=\mathcal A_{\mathrm{KG}}\,,\qquad \mathcal A_I=\frac{\lambda}{4!}\int\varphi^4(x)\dd^4x
\]
(later we will see that in actual calculations it is more convenient to include
certain parts of the quadratic action, the so called counterterms, into
$\mathcal A_I$, but at the moment we ignore this subtlety). With this
splitting, the power series expansion of $\exp(-\mathcal A_I)$ generates
expansion in the powers of $\lambda$.

Suppose we want to calculate such expansion of the two-point correlation
function $\langle\varphi(x_1)\varphi(x_2)\rangle$. By its definition
\[
\bigl\langle\varphi(x_1)\varphi(x_2)\bigr\rangle=Z^{-1}\int\bigl[D\varphi\bigr]\,\varphi(x_1)\varphi(x_2)\,e^{-\mathcal A_0}\,e^{-\mathcal A_I}\,.
\]
Expanding the last exponential as prescribed, we have
\[
Z^{-1}\int\bigl[D\varphi\bigr]\,\varphi(x_1)\varphi(x_2)\left(1-\frac{\lambda}{4!}\int\varphi^4(x)\dd^4x+O(\lambda^2)\right)\,e^{-\mathcal A_0}\,.
\]
The zero-order term is just the functional integral for the two-point function
of the KG theory. We can rewrite the previous expression as
\begin{equation}\label{eq:pt2}
\bigl\langle\varphi(x_1)\varphi(x_2)\bigr\rangle=\frac{Z_0}{Z}\left[\bigl\langle\varphi(x_1)\varphi(x_2)\bigr\rangle_0
-\frac{\lambda}{4!}\int\bigl\langle\varphi(x_1)\varphi(x_2)\varphi^4(x)\bigr\rangle_0\dd^4x+\dots\right]\,,
\end{equation}
where $Z_0=\int D[\varphi]\,e^{-\mathcal A_0}$ and $\langle\cdots\rangle_0$
stand for the partition function and correlation functions in the free theory.
This expression gives an example of the general situation: the perturbation
theory allows one to write the correlation functions of interacting theory as a
power series in the coupling parameter, the coefficients being certain
integrals of the correlation functions of the free theory. The role of the
factor $\frac{Z_0}{Z}$ will be clarified later.

The correlation functions of the free theory are computed using the Wick's
rules. In the above example, the first-order in $\lambda$ term involves the
$\dd^4x$ integral of the correlation function
\[
\bigl\langle\varphi(x_1)\varphi(x_2)\,\varphi^4(x)\bigr\rangle_0=\bigl\langle\varphi(x_1)\varphi(x_2)\varphi(x)\varphi(x)\varphi(x)\varphi(x)\bigr\rangle_0\,.
\]
This is a six-point correlation function of the free theory, with four points
coincident at $x$. According to the Wick's rule, it can be written as the sum
of
\[
5!!=15
\]
terms corresponding to different ways of pairwise contractions of the six field
insertions involved. Many different contractions lead to identical
contributions. There are two classes of such contractions with essentially
different contributions:
\begin{enumerate}[label=(\roman*)]
\item First, one can contract $\varphi(x_1)\varphi(x_2)\to
\contraction{}{\varphi(x_1)}{}{\varphi(x_2)}\varphi(x_1)\varphi(x_2)$, and then
contract four factors $\varphi(x)$ among themselves. There are three
possibilities for contracting. Say, one first contracts the first factor in
$\varphi(x)\varphi(x)\varphi(x)\varphi(x)$ with any one of the other three (3
ways), and once that one is chosen, there is only one way to contract the
remaining two insertions. We find that such contractions bring in the
contribution
\[
3\times
\contraction{}{\varphi(x_1)}{}{\varphi(x_2)}
\contraction{\varphi(x_1)\varphi(x_2)\,}{\varphi(x)}{}{\varphi(x)}
\contraction{\varphi(x_1)\varphi(x_2)\,\varphi(x)\varphi(x)\,}{\varphi(x)}{}{\varphi(x)}
\varphi(x_1)\varphi(x_2)\,\varphi(x)\varphi(x)\,\varphi(x)\varphi(x)
=3\,D(x_1-x_2)\,D(x-x)\,D(x-x)\,.
\]

\item Alternatively, one can contract $\varphi(x_1)$ with one of the four
identical factors $\varphi(x)$ (4 ways), and $\varphi(x_2)$ with any one of the
remaining three factors $\varphi(x)$ (3 ways); this leaves just two insertions
$\varphi(x)$, with only one way of contractions. Therefore there are
$4\times3=12$ identical contributions
\[
12\times
\contraction[2ex]{}{\varphi(x_1)}{\varphi(x_2)}{\varphi(x)}
\contraction{\varphi(x_1)}{\varphi(x_2)}{\varphi(x)}{\varphi(x)}
\contraction{\varphi(x_1)\varphi(x_2)\varphi(x)\varphi(x)}{\varphi(x)}{}{\varphi(x)}
\varphi(x_1)\varphi(x_2)\varphi(x)\varphi(x)\varphi(x)\varphi(x)
=12\,D(x_1-x)\,D(x_2-x)\,D(x-x)\,.
\]
\end{enumerate}
The resulting contribution $\sim\lambda$ in the brackets in \eqref{eq:pt2}
then reads
\begin{equation}\label{eq:pt2first}
-\frac{\lambda}{4!}\left[3\,D(x_1-x_2)\int \dd^4x\,D^2(x-x)+12\int \dd^4x\,D(x_1-x)\,D(x_2-x)\,D(x-x)\right]\,.
\end{equation}
Depicting contractions as solid lines
\[
\contraction{}{\varphi(x)}{}{\varphi(x')}\varphi(x)\varphi(x')
=\ x\;\dia{\pic{line};}\;x'\ \to\ D(x-x')
\]
the two terms are represented by two diagrams
\begin{equation}\label{eq:pt2diag}
3\times\Bigl[\,\dia{\pic{line}; \lab[left]{-1.4,0}{x_1} \lab[right]{1.4,0}{x_2} \pic at (3.4,0) {eight};}\,\Bigr]
+12\times\Bigl[\,\dia{\pic{tadline}; \lab[left]{-1.4,0}{x_1} \lab[right]{1.4,0}{x_2}}\,\Bigr]
\end{equation}
where the loop represents the contraction
\[
\contraction{}{\varphi(x)}{}{\varphi(x)}\varphi(x)\varphi(x)
=\dia{\draw[prop] (0,0.42) circle (0.42); \node[vtx] at (0,0) {};}\ \to\ D(x-x)\,.
\]
Recalling the overall factor
\[
-\frac{\lambda}{4!}=-\frac{\lambda}{24}
\]
in \eqref{eq:pt2}, we find that the two diagrams in \eqref{eq:pt2diag} come with
the factors $-\lambda/S_1$ and $-\lambda/S_2$, respectively, where the so
called ``symmetry factors'' here are
\[
\frac1{S_1}=\frac18\,,\qquad\frac1{S_2}=\frac12\,.
\]

Similar structure is found in the higher orders of the $\lambda$-expansion. The
$n$-th order term has the form
\begin{equation}\label{eq:ptn}
\frac1{n!}\left(-\frac{\lambda}{4!}\right)^n\left\langle\varphi(x_1)\varphi(x_2)\left[\int\varphi^4(y)\dd^4y\right]^n\right\rangle_0
\end{equation}
which involves the $n$-fold 4-space integration
\[
\left[\int\varphi^4(y)\dd^4y\right]^n=\int\varphi^4(y_1)\varphi^4(y_2)\cdots\varphi^4(y_n)\dd^4y_1\cdots \dd^4y_n\,.
\]
The correlation function in the integrand is that of the product of $2+4n$
factors $\varphi$, two at $x_1$ and $x_2$, and $4n$ at the integration points
$y_1,\dots,y_n$, four factors at each point.

Contractions of these field factors are represented by diagrams. The diagrams
are built from $2n+1$ contraction lines connecting two types of ``vertices'':

There are two ``external vertices'' associated with the points $x_1$ and
$x_2$, with exactly one of the contraction lines attached to each,
\[
\text{External points:}\qquad x_1\ \dia{\draw[prop] (0,0) node[vtx]{} -- (1.4,0);}
\qquad\qquad \dia{\draw[prop] (0,0) -- (1.4,0) node[vtx]{};}\ x_2
\]
and $n$ ``interaction vertices'', associated with $n$ integration points
$y_1,y_2,\dots,y_n$, with four contraction lines connected to each of them:
\[
\text{Interaction points:}\qquad
\dia{\pic{xvertex}; \lab[right]{0.3,-0.05}{y_1}}\quad
\dia{\pic{xvertex}; \lab[right]{0.3,-0.05}{y_2}}\quad\cdots\quad
\dia{\pic{xvertex}; \lab[right]{0.3,-0.05}{y_n}}
\]
Thus, in the order $\lambda^2$ the diagrams involve two external and two
interaction vertices, for example
\[
x_1\ \dia{\pic{sunline};}\ x_2
\]
Again, many distinct contractions are represented by identical diagrams. For
instance, the above diagram may represent the contraction
\[
\contraction{}{\varphi(x_1)}{}{\varphi(y_1)}
\contraction[4ex]{\varphi(x_1)\varphi(y_1)}{\varphi(y_1)}{\varphi(y_1)\varphi(y_1)\varphi(y_2)\varphi(y_2)}{\varphi(y_2)}
\contraction[3ex]{\varphi(x_1)\varphi(y_1)\varphi(y_1)}{\varphi(y_1)}{\varphi(y_1)\varphi(y_2)}{\varphi(y_2)}
\contraction[2ex]{\varphi(x_1)\varphi(y_1)\varphi(y_1)\varphi(y_1)}{\varphi(y_1)}{}{\varphi(y_2)}
\contraction{\varphi(x_1)\varphi(y_1)\varphi(y_1)\varphi(y_1)\varphi(y_1)\varphi(y_2)\varphi(y_2)\varphi(y_2)}{\varphi(y_2)}{}{\varphi(x_2)}
\varphi(x_1)\varphi(y_1)\varphi(y_1)\varphi(y_1)\varphi(y_1)\varphi(y_2)\varphi(y_2)\varphi(y_2)\varphi(y_2)\varphi(x_2)
\]
as well as any one of
\[
4\times4\times3!
\]
equivalent contractions related to this one by interchanging of the contracted
factors $\varphi$ inside $\varphi^4(y_1)$ and $\varphi^4(y_2)$. Every such
contraction corresponds to the same product of the propagators
\[
D(x_1-y_1)\,D^3(y_1-y_2)\,D(y_2-x_2)\,,
\]
and since the contributions involve the integrations over the interaction
points
\[
\int \dd^4y_1\dd^4y_2\,(\dots)
\]
the diagrams related by the interchange $y_1\leftrightarrow y_2$ yield
identical contributions. Therefore the above contraction comes with the
multiplicity factor
\[
2\cdot4\cdot4!\,.
\]
This factor almost entirely cancels the denominators in the overall factors
$(1/n!)(-\lambda/4!)^n$ in \eqref{eq:ptn}, which for $n=2$ is
\[
\frac1{2!}\left(-\frac{\lambda}{4!}\right)^2\,,
\]
leaving behind the factor $(-\lambda)^2/3!$. Finally, total contribution of
the contractions corresponding to the diagram in this example is
\[
\frac{(-\lambda)^2}{S}\int \dd^4y_1\dd^4y_2\,D(x_1-y_1)\,D^3(y_1-y_2)\,D(y_2-x_2)\,,
\]
with the ``symmetry factor'' $\frac1S=\frac1{3!}$.

This example illustrates the general rules of the diagrammatic representations
of the contributions in the perturbation theory, known as the

\subsection*{Feynman Rules (for \texorpdfstring{$\varphi^4$}{phi**4})}

To find $\lambda^n$ contribution to
\[
\bigl\langle\varphi(x_1)\cdots\varphi(x_N)\bigr\rangle
\]
\begin{enumerate}
\item Draw all diagrams with $N$ external vertices associated with the points
$x_1,\dots,x_N$, and with $n$ interaction vertices. Exclude disconnected
diagrams with ``vacuum fragments'' (see below).
\item For each diagram, associate
\begin{enumerate}[label=(\roman*)]
\item $x\ \dia{\pic{line};}\ x'=D(x-x')$;
\item $\dia{\pic{xvertex}; \lab[right]{0.3,-0.05}{y}}=-\lambda\int \dd^4y$;
\item Multiply by the ``symmetry factor'' $1/S$.
\end{enumerate}
\item Sum the contributions of all the diagrams.
\end{enumerate}

\begin{remark}
These rules are generalized in a straightforward way to interactions involving
any powers of the field $\varphi$ (possibly including terms linear and
quadratic in $\varphi$)
\begin{equation}\label{eq:genint}
\mathcal A_I=\int \dd^4y\,\left[\lambda_1(y)\,\varphi(y)+\frac{\lambda_2(y)}{2!}\,\varphi^2(y)+\frac{\lambda_3(y)}{3!}\,\varphi^3(y)+\frac{\lambda_4(y)}{4!}\,\varphi^4(y)+\cdots\right]\,.
\end{equation}
Here I allowed the coupling parameters $\lambda_k=\lambda_k(y)$ to be
arbitrary functions of $y$. (This generally breaks space-time translation and
rotation symmetries of the theory, but otherwise defines quite meaningful
quantity. Later, these generalizations will allow us to use diagrammatic
representations of the generating functional and effective action.) The
Feynman rules generalize to the case \eqref{eq:genint} by including $k$-leg
vertices
\begin{equation}\label{eq:kleg}
\begin{gathered}
\dia{\draw[prop] (-1.4,0) -- (0,0) node[vtx]{};}=-\int \dd^4y\,\lambda_1(y)\dots\,,\qquad
\dia{\draw[prop] (-1.4,0) -- (1.4,0); \node[vtx] at (0,0) {};}=-\int \dd^4y\,\lambda_2(y)\dots\,,\\[4pt]
\dia{\draw[prop] (-1,0.8) -- (0,0) -- (-1,-0.8); \draw[prop] (0,0) -- (1.3,0); \node[vtx] at (0,0) {};}=-\int \dd^4y\,\lambda_3(y)\dots\,,\qquad\cdots\,,\qquad
\dia{\foreach \a in {20,60,100,140,250,290} \draw[prop] (0,0) -- (\a:1.1);
     \draw[densely dotted] (165:0.7) arc[start angle=165,end angle=230,radius=0.7];
     \node[vtx] at (0,0) {}; \lab[below]{0.25,-0.9}{k}}=-\int \dd^4y\,\lambda_k(y)\dots\,.
\end{gathered}
\end{equation}
\end{remark}

\subsection*{Vacuum Diagrams}

Returning back to the homogeneous $\varphi^4$ theory, let us take another look
at the diagrams contributing to the two-point function. We write
\begin{equation}\label{eq:twofactors}
\bigl\langle\varphi(x_1)\varphi(x_2)\bigr\rangle
=\left[\frac1{Z_0}\int\bigl[D\varphi\bigr]\,\varphi(x_1)\varphi(x_2)\,e^{-\mathcal A_0-\mathcal A_I}\right]\times\left[\frac{Z_0}{Z}\right]\,,
\end{equation}
where $Z_0$ is the partition function of the free theory defined by
$\mathcal A_0=\mathcal A_{\mathrm{KG}}$. The diagrams we wrote down above
actually correspond to the first factor in \eqref{eq:twofactors}. For this
first factor we have
\[
\underbrace{\dia{\pic{line}; \lab[left]{-1.4,0}{x_1} \lab[right]{1.4,0}{x_2}}}_{(a)}
+\frac12\,\Bigl(\underbrace{\dia{\pic{tadline}; \lab[left]{-1.4,0}{x_1} \lab[right]{1.4,0}{x_2}}}_{(b)}\Bigr)
+\frac18\,\Bigl(\underbrace{\dia{\pic{line}; \lab[left]{-1.4,0}{x_1} \lab[right]{1.4,0}{x_2} \pic at (3.3,0) {eight};}}_{(c)}\Bigr)+\cdots\,.
\]
Observe that there are diagrams containing disconnected pieces, like the
diagram (c) above. This diagram has a fragment
\[
\dia{\pic{eight};}
\]
which is not connected to any of the ``external'' points ($x_1,x_2$ in this
case). Such fragments are called ``vacuum diagrams'' (or ``vacuum
fragments''). The value
\[
D(x_1-x_2)\times\left(-\frac{\lambda}{8}\int \dd^4x\,D(x-x)\,D(x-x)\right)
\]
associated with the diagram (c) above is the product of the values of its
parts,
\begin{equation}\label{eq:eightvalue}
\dia{\pic{line}; \lab[left]{-1.4,0}{x_1} \lab[right]{1.4,0}{x_2}}=D(x_1-x_2)\,;\qquad
\dia{\pic{eight};}=-\frac{\lambda}{8}\int \dd^4x\,D^2(x-x)\,.
\end{equation}
As follows directly from the Feynman rules above, this holds in general: the
value of any disconnected diagram is the product of the values of its
connected fragments. In fact, in computing the correlation function
\eqref{eq:twofactors} the vacuum fragments should be ignored because their
contributions are canceled exactly by the second factor in the r.h.s.\ of
\eqref{eq:twofactors}, which contains the partition function $Z$ of the full
theory. To see this, consider again the diagram (c) above. In higher orders of
the perturbation theory we obtain infinitely many diagrams in which the vacuum
fragment in (c) is repeated many times,
\[
\Bigl(\dia{\pic{line}; \lab[left]{-1.4,0}{x_1} \lab[right]{1.4,0}{x_2}}\quad
\underbrace{\dia{\pic{eight}; \pic at (1.3,0) {eight}; \node at (2.25,0) {$\cdots$}; \pic at (3.2,0) {eight};}}_{n\ \text{times}}\Bigr)
\]
The value corresponding to $n$ copies of the vacuum fragment is just
\[
\underbrace{\dia{\pic{eight}; \pic at (1.3,0) {eight}; \node at (2.25,0) {$\cdots$}; \pic at (3.2,0) {eight};}}_{n\ \text{times}}
=\frac1{n!}\,\Bigl[\dia{\pic{eight};}\Bigr]^n\,,
\]
the factor $1/n!$ being the extra symmetry factor associated with the
permutations of the fragments. The sum over $n$ therefore yields
\[
\Bigl(\dia{\pic{line}; \lab[left]{-1.4,0}{x_1} \lab[right]{1.4,0}{x_2}}\Bigr)\times\exp\Bigl(\dia{\pic{eight};}\Bigr)
\]
This simple analysis can be repeated with any other vacuum fragment. As the
result we obtain for the first factor in \eqref{eq:twofactors}
\begin{multline*}
\Bigl(\underbrace{\dia{\pic{line}; \lab[left]{-1.4,0}{x_1} \lab[right]{1.4,0}{x_2}}
+\dia{\pic{tadline}; \lab[left]{-1.4,0}{x_1} \lab[right]{1.4,0}{x_2}}
+\dia{\pic{sunline}; \lab[left]{-1.7,0}{x_1} \lab[right]{1.7,0}{x_2}}+\cdots}_{\text{sum of all connected diagrams}}\Bigr)\\
\times\exp\Bigl(\underbrace{\dia{\pic{eight};}+\dia{\pic{melon};}+\dia{\pic{chain3};}+\cdots}_{\mathclap{\text{sum of all connected vacuum diagrams}}}\Bigr)
\end{multline*}
Consider now the second factor in \eqref{eq:twofactors},
\[
\frac Z{Z_0}=\frac1{Z_0}\int\bigl[D\varphi\bigr]\,e^{-\mathcal A_0}\,e^{-\mathcal A_I}\,.
\]
Expanding in $\mathcal A_I$ we find that this quantity is given by the sum of
all vacuum diagrams, which again exponentiates in terms of the connected
vacuum diagrams. That is
\begin{equation}\label{eq:ZZ0}
Z/Z_0=\exp\Bigl(\text{sum of all connected vacuum diagrams}\Bigr)
\end{equation}
so that the second factor in \eqref{eq:twofactors} cancels the contributions of
all vacuum fragments. Thus, if we are interested in the correlation function,
the factor $Z^{-1}$ in the definition of the correlation function can indeed be
ignored if we simultaneously drop all the diagrams with disconnected vacuum
fragments.

Note that the above analysis of the ``vacuum fragments'' generalizes
straightforwardly to the case of interaction \eqref{eq:genint} with arbitrary
$x$-dependent coupling constants. In particular it is still true that the
partition function of this more general theory exponentiates in terms of the
connected vacuum diagrams, as in \eqref{eq:ZZ0}. Of course in this case the
diagrams involve all the vertices in \eqref{eq:kleg}.

Let me comment on the interpretation of the above formula \eqref{eq:ZZ0} for
$Z$. The first of the connected vacuum diagrams is given by the vacuum
fragment in \eqref{eq:eightvalue}, i.e.
\[
\dia{\pic{eight};}=-\frac{\lambda}{8}\,D^2(0)\int \dd^4x
\]
Clearly, it contains an infinite factor
\[
\int \dd^4x=V^{(4)}\,,
\]
where $V^{(4)}$ is the volume of the 4-dimensional space. This is the common
feature of the connected vacuum diagrams
\[
\dia{\pic{eight};}\,,\qquad\dia{\pic{melon};}\,,\qquad\dia{\pic{chain3};}\,,\qquad\dots
\]
As a consequence of translational invariance each of them contains this volume
factor. Therefore
\begin{equation}\label{eq:Zvac}
Z=Z_0\,\exp\bigl(-\varepsilon_I\,V^{(4)}\bigr)
\end{equation}
where $\varepsilon_I$ is just the sum of all connected vacuum diagrams with the
volume factored out (and with the sign reversed).

The functional integral
\[
Z=\int\bigl[D\varphi\bigr]\,e^{-\mathcal A}
\]
can be interpreted as a (configuration) partition function of a classical
statistical mechanics. The classical system in this interpretation is a field
$\varphi(x)$ in 4-dimensional (Euclidean) space, with the potential energy
$U[\varphi]$,
\[
\beta\,U[\varphi]:=\frac1\hbar\,\mathcal A[\varphi]\,,
\]
where $\beta$ is interpreted as the inverse temperature. ($\beta$ can be set
equal to 1, again by suitable choice of the energy units.) This is a spatially
homogeneous system with infinitely many degrees of freedom associated with all
points in the 4-dimensional space. Naturally, its free energy is an intensive
quantity proportional to $V^{(4)}$. We see that the quantity
\[
\varepsilon_0+\varepsilon_I\,,\qquad\text{where}\qquad \varepsilon_0=-\log Z_0/V^{(4)}
\]
is interpreted as the specific (per unit volume) free energy of this system,
the first term being the specific free energy of the free theory and
$\varepsilon_I$ incorporating all corrections due to the interaction.

Another way to look at the formula \eqref{eq:ZZ0} is to recall the relation
between the path integrals and the matrix elements of the (imaginary) time
evolution operator. The path integral defining $Z$ can be thought of as the
$T\to\infty$ limit of the path integral
\[
\int\bigl[D\varphi\bigr]_{\substack{\varphi(\xx,-T/2)=\varphi_i(\xx)\\\varphi(\xx,+T/2)=\varphi_f(\xx)}}\,e^{-\mathcal A}
\]
which represents the matrix element
\[
\bra{\varphi_f(\xx)}e^{-T\hat H}\ket{\varphi_i(\xx)}\,.
\]
Here $\varphi_i$ and $\varphi_f$ are arbitrary ``initial'' and ``final''
field configurations. The evolution operator inside can be written as
\[
e^{-T\hat H}=\sum_n\ket n\bra n\,e^{-TE_n}\,,
\]
where $\ket n$ is a complete set of the stationary states; its $T\to\infty$
limit then is
\[
e^{-E_0T}\,\ket0\bra0\,,\qquad\text{where}\qquad E_0=\varepsilon\,V^{(3)}
\]
($V^{(3)}$ is the volume of the 3-dimensional ``equal time'' section of our
4-dimensional space). Because $V^{(4)}=T\times V^{(3)}$, comparison with the
above formula \eqref{eq:ZZ0} shows that vacuum diagrams represent the
expansion of the ground state energy in the powers of $\lambda$ (hence the term
``vacuum diagrams''). We will see later that all the vacuum diagrams are highly
divergent at short distances. This is similar to the short-distance (high
momenta) divergence of the vacuum energy of the Klein-Gordon theory which we
observed before. Proper understanding of the physical content of these
divergent terms is important in certain problems like the Casimir effect or
interaction with gravity.
